% ================= CAPÍTULO 1: CONJUNTOS =================
\chapter{Conjuntos}
\prob{SEMANA 1}{la base del lenguaje: todo lo demás se escribe con conjuntos}

\section{Inclusión e igualdad}
\agregado
\begin{sello}{DEFINICIONES}
$A\subset B$ ($A$ está incluido en $B$) $\iff$ todo elemento de $A$ es elemento de $B$.\\[3pt]
$A=B\iff A\subset B$ y $B\subset A$. \quad $\emptyset$ es el conjunto sin elementos; $\emptyset\subset A$ para todo $A$.
\end{sello}
\begin{metodo}{PROBAR UNA INCLUSIÓN}
\begin{enumerate}
\item Escribí «sea $x\in A$».
\item Usá las definiciones para llegar a «entonces $x\in B$».
\item Para una igualdad, hacelo en los dos sentidos.
\end{enumerate}
\end{metodo}

\section{Unión}
\apuntes
\begin{sello}{DEFINICIÓN Y PROPIEDADES}
$A\cup B=\{x: x\in A\ \text{o}\ x\in B\}$.\\[3pt]
1) $A\cup A=A$ \quad 2) $A\cup B=B\cup A$ \quad 3) $A\cup\emptyset=A$ \quad 4) $A\subset A\cup B$ \quad 5) $B\subset A\iff B\cup A=A$
\end{sello}
\begin{demo}
\textbf{Prop.\ 3.} $A\cup\emptyset=\{x:x\in A\text{ o }x\in\emptyset\}=\{x:x\in A\}=A$, porque «$x\in\emptyset$» nunca pasa.\\[3pt]
\textbf{Prop.\ 5 ($\Rightarrow$).} Supongo $B\subset A$. Quiero $B\cup A=A$: dos inclusiones.
$A\subset B\cup A$ por la prop.\ 4. Para $B\cup A\subset A$: sea $x\in B\cup A$; entonces $x\in B$ o $x\in A$. Si $x\in B$, como $B\subset A$, $x\in A$. Si no, ya $x\in A$. En los dos casos $x\in A$.\\[3pt]
\textbf{Prop.\ 5 ($\Leftarrow$).} Supongo $B\cup A=A$. Sea $x\in B$: entonces $x\in B\cup A=A$. Luego $B\subset A$. $\blacksquare$
\end{demo}
\begin{tip}{LA IMAGEN}
Pegarle $B$ a $A$ no cambia nada exactamente cuando $B$ ya estaba adentro de $A$.
\end{tip}

\section{Intersección}
\apuntes
\begin{sello}{DEFINICIÓN Y PROPIEDADES}
$A\cap B=\{x: x\in A\ \text{y}\ x\in B\}$. \ Ej.: $\{1,2,5\}\cap\{1,5,7,21\}=\{1,5\}$.\\[3pt]
1) $A\cap A=A$ \quad 2) $A\cap B=B\cap A$ \quad 3) $A\cap\emptyset=\emptyset$ \quad 4) $A\cap B\subset A$ \quad 5) $B\subset A\iff B\cap A=B$
\end{sello}
\begin{demo}
\textbf{Prop.\ 4.} Sea $x\in A\cap B$: por definición $x\in A$.\\[3pt]
\textbf{Prop.\ 5 ($\Leftarrow$).} $B=B\cap A\subset A$ (props.\ 4 y 2).\\[3pt]
\textbf{Prop.\ 5 ($\Rightarrow$).} $B\cap A\subset B$ por la prop.\ 4. Para $B\subset B\cap A$: sea $x\in B$; como $B\subset A$, también $x\in A$; luego $x\in B\cap A$. $\blacksquare$
\end{demo}

\agregado
\begin{sello}{ASOCIATIVAS Y DISTRIBUTIVAS}
$(A\cup B)\cup C=A\cup(B\cup C)$, \quad $(A\cap B)\cap C=A\cap(B\cap C)$.\\[3pt]
$A\cap(B\cup C)=(A\cap B)\cup(A\cap C)$, \quad $A\cup(B\cap C)=(A\cup B)\cap(A\cup C)$.
\end{sello}
\begin{demo}
Primera distributiva. $x\in A\cap(B\cup C)\iff x\in A$ y ($x\in B$ o $x\in C$) $\iff$ ($x\in A$ y $x\in B$) o ($x\in A$ y $x\in C$) $\iff x\in(A\cap B)\cup(A\cap C)$. La otra es igual cambiando «y» por «o». $\blacksquare$
\end{demo}

\section{Diferencia}
\apuntes
\begin{sello}{DEFINICIÓN Y PROPIEDADES}
$A\setminus B=\{x\in A: x\notin B\}$. \ Ej.: $\{1,2,3\}\setminus\{1,3,21,23\}=\{2\}$. En general $A\setminus B\neq B\setminus A$.\\[3pt]
1) $A\setminus A=\emptyset$ \quad 2) $A\setminus\emptyset=A$ \quad 3) $\emptyset\setminus A=\emptyset$ \quad 4) $A\setminus(B\cup C)=(A\setminus B)\setminus C$
\end{sello}
\agregado
\begin{demo}
\textbf{Prop.\ 4.} $x\in A\setminus(B\cup C)\iff x\in A$ y $x\notin B\cup C\iff x\in A$, $x\notin B$ y $x\notin C\iff x\in(A\setminus B)$ y $x\notin C\iff x\in(A\setminus B)\setminus C$. $\blacksquare$
\end{demo}
\begin{sello}{EJERCICIO DEL CUADERNO: $A\setminus B=B\setminus A\iff A=B$}
($\Leftarrow$) Si $A=B$, los dos lados son $A\setminus A=\emptyset$.\\[2pt]
($\Rightarrow$) Si hubiera $x\in A\setminus B$, estaría en $B\setminus A$, o sea $x\in B$; pero $x\notin B$. Absurdo: $A\setminus B=\emptyset$, es decir $A\subset B$. Igual $B\setminus A=\emptyset$ da $B\subset A$. Luego $A=B$. $\blacksquare$
\end{sello}

\section{Complemento y De Morgan}
\apuntes
\begin{sello}{COMPLEMENTO}
Dado un universo $U$ y $A\subset U$: $A^c=\{x\in U:x\notin A\}$.\\[3pt]
Observación: $A\setminus B=A\cap B^c$. \quad $\emptyset^c=U$, \ $U^c=\emptyset$.
\end{sello}
\begin{sello}{LEYES DE DE MORGAN}
$(A\cap B)^c=A^c\cup B^c$ \qquad $(A\cup B)^c=A^c\cap B^c$
\end{sello}
\agregado
\begin{demo}
$x\in(A\cap B)^c\iff$ no ($x\in A$ y $x\in B$) $\iff x\notin A$ o $x\notin B\iff x\in A^c\cup B^c$. La segunda es igual. $\blacksquare$
\end{demo}
\begin{tip}{EN UNA FRASE}
El complemento da vuelta todo: «y» pasa a «o», «$\cap$» pasa a «$\cup$». Es la negación de la lógica (\emph{no}(P y Q) = \emph{no} P o \emph{no} Q) escrita con conjuntos.
\end{tip}

\section{Producto cartesiano}
\agregado
\begin{sello}{PARES ORDENADOS}
$A\times B=\{(a,b): a\in A,\ b\in B\}$. \ $(a,b)=(a',b')\iff a=a'$ y $b=b'$. El orden importa.\\[2pt]
$\R\times\R=\R^2$ es el plano: cada punto es un par $(x,y)$.
\end{sello}
